\documentclass{article}

% NeurIPS 2026 style (use neurips_2026.sty if available; fallback to standard)
\usepackage[final]{neurips_2026}
\usepackage[utf8]{inputenc}
\usepackage[T1]{fontenc}
\usepackage{hyperref}
\usepackage{url}
\usepackage{booktabs}
\usepackage{amsfonts}
\usepackage{amsmath}
\usepackage{amssymb}
\usepackage{nicefrac}
\usepackage{microtype}
\usepackage{xcolor}
\usepackage{graphicx}
\usepackage{multirow}
\usepackage{subcaption}
\usepackage{algorithm}
\usepackage{algorithmic}
\usepackage{tablefootnote}

\title{Structured Social State for NPC Dialogue Generation:\\Architecture Benchmarks and Latent Conditioning}

\author{%
  Anonymous Authors\\
  \texttt{anonymous@institution.edu}
}

\begin{document}

\maketitle

% ═══════════════════════════════════════════════════════════════════════════════
% ABSTRACT
% ═══════════════════════════════════════════════════════════════════════════════
\begin{abstract}
We present a framework for social-state-conditioned NPC dialogue generation that explicitly models 29 latent social dimensions spanning conversational intent, affect, theory of mind, relational stance, normative pressures, and response policy.
We benchmark four from-scratch small language model (SLM) architectures (GPT, PrefixGPT, MoE, Mamba-like; 15--22M parameters), two conditioning encoders (personality OCEAN, affect VAD), four dialogue response models, and a structured 29-head latent state predictor built on Qwen3-1.7B.
Key findings: (1)~Mixture-of-Experts achieves the best from-scratch perplexity ($\text{PPL}=42.07$, $-7.2\%$ vs.\ standard GPT); (2)~explicit OCEAN+VAD soft-prefix conditioning reduces response perplexity by $12.3\%$; (3)~the 29-head predictor attains mean accuracy $0.702$ (estimated $\kappa=0.611$) across evaluated heads, with response policy $\text{F1}$ in the $0.45$--$0.51$ range across evaluation artifacts; and (4)~keyword-based automatic checks detect zero secret leakage in sampled generations.
Our structured latent state $Z_t$ provides an interpretable, auditable intermediate representation that bridges perception and generation for interactive NPC dialogue.
\end{abstract}

% ═══════════════════════════════════════════════════════════════════════════════
% 1. INTRODUCTION
% ═══════════════════════════════════════════════════════════════════════════════
\section{Introduction}
\label{sec:intro}

Non-player characters (NPCs) in interactive narratives must exhibit consistent social behaviour: remembering relationships, keeping secrets, reacting emotionally, and selecting contextually appropriate responses.
Current approaches either rely on hand-scripted dialogue trees---which are rigid and expensive to author---or fine-tune large language models (LLMs) without explicit social-state representations, yielding fluent but socially inconsistent outputs~\cite{park2023generative,wang2023voyager}.

We propose a \textbf{structured latent social state} $Z_t$ that decomposes each dialogue turn into 29 classification targets organised into six groups:
\begin{itemize}
    \item \textbf{$C_t$} (Conversational): dialogue act, tone, risk type
    \item \textbf{$A_t$} (Affect): valence, arousal, threat, control
    \item \textbf{$M_t$} (Mental model): player intent, knowledge, credibility
    \item \textbf{$R_t$} (Relational stance): 6 dimensions $\times$ level + delta
    \item \textbf{$N_t$} (Normative pressure): duty, secrecy, face, value conflict
    \item \textbf{$D_t$} (Decision policy): response policy, reveal decision, repair strategy
\end{itemize}

This decomposition is motivated by sociological and pragmatic frameworks~\cite{brown1987politeness,oatley1987towards,spencer2008politeness} and enables interpretable, auditable NPC behaviour.
Our research questions are:

\begin{description}
    \item[RQ1] Which from-scratch SLM architecture performs best on NPC dialogue at the 15--22M parameter scale?
    \item[RQ2] Can a pretrained LLM learn to predict the structured latent social state from dialogue context alone?
    \item[RQ3] Does explicit social-state conditioning improve response generation quality?
\end{description}

\paragraph{Contributions.}
(1)~A 29-dimensional structured social-state schema for NPC dialogue with formal consistency rules;
(2)~a four-track benchmark spanning 13 trained or exploratory models from 15M to 4B parameters;
(3)~empirical evidence that explicit social-state conditioning reduces perplexity by $12.3\%$;
(4)~a multi-head latent predictor achieving mean estimated $\kappa=0.611$ across evaluated heads at the 1.7B scale.

% ═══════════════════════════════════════════════════════════════════════════════
% 2. RELATED WORK
% ═══════════════════════════════════════════════════════════════════════════════
\section{Related Work}
\label{sec:related}

\paragraph{NPC dialogue.}
Prior work on NPC dialogue systems ranges from finite-state machines and dialogue trees~\cite{walker2006individual} to neural approaches using sequence-to-sequence models~\cite{urbanek2019learning}.
Recent work has explored LLM-based agents for open-world games~\cite{park2023generative,wang2023voyager}, but without structured social-state modelling.

\paragraph{Social simulation.}
Computational models of social interaction draw on affect theory~\cite{oatley1987towards,russell1980circumplex}, personality modelling (Big Five / OCEAN)~\cite{digman1990personality}, and politeness theory~\cite{brown1987politeness}.
Our framework integrates these into a unified latent state that is both predictable and actionable.

\paragraph{Conditioned language generation.}
Prefix tuning~\cite{li2021prefix}, control codes~\cite{keskar2019ctrl}, and soft prompting~\cite{lester2021power} enable attribute-conditioned text generation.
We extend soft-prefix conditioning to a multi-dimensional social state (OCEAN + VAD) and demonstrate its effectiveness for NPC dialogue.

\paragraph{Multi-task classification on LLMs.}
Multi-head classifiers on top of pretrained LLMs have been explored for various NLP tasks~\cite{liu2019multi}.
Our 29-head predictor is, to our knowledge, the first application to structured social-state prediction for interactive dialogue.

% ═══════════════════════════════════════════════════════════════════════════════
% 3. METHOD
% ═══════════════════════════════════════════════════════════════════════════════
\section{Method}
\label{sec:method}

\subsection{Social State Schema ($Z_t$)}
\label{sec:schema}

Each dialogue turn $t$ is annotated with a structured social state $Z_t = (C_t, A_t, M_t, R_t, N_t, D_t)$, comprising 29 classification targets.
Table~\ref{tab:schema} summarises the schema.

\begin{table}[t]
\centering
\caption{Structured social state schema $Z_t$. Each field is a classification target predicted by the 29-head model. Stance dimensions ($R_t$) each have a level (5-class ordinal) and a per-turn delta (5-class ordinal).}
\label{tab:schema}
\small
\begin{tabular}{@{}llrl@{}}
\toprule
\textbf{Group} & \textbf{Field} & \textbf{\#Classes} & \textbf{Type} \\
\midrule
$C_t$ & dialogue\_act & 10 & multi-label \\
      & tone & 6 & single-label \\
      & risk\_type & 5 & single-label \\
\midrule
$A_t$ & valence, arousal, threat, control & 3 each & single-label \\
\midrule
$M_t$ & player\_intent & 9 & single-label \\
      & player\_knowledge & 4 & single-label \\
      & player\_credibility & 3 & single-label \\
\midrule
$R_t$ & \{affection, respect, dominance, & 5 each & ordinal \\
      & \ familiarity, trust, obligation\}$_\text{level}$ & & \\
      & \{...\}$_\text{delta}$ & 5 each & ordinal \\
\midrule
$N_t$ & duty, secrecy, face pressure & 3 each & single-label \\
      & value\_conflict & 3 & single-label \\
\midrule
$D_t$ & response\_policy & 10 & single-label \\
      & reveal\_decision & 4 & single-label \\
      & repair\_strategy & 5 & single-label \\
\bottomrule
\end{tabular}
\end{table}

\paragraph{Consistency rules.}
We enforce domain constraints: e.g., $\text{reveal\_decision} = \textit{full}$ requires $\text{secrecy\_pressure} \neq \textit{high}$; $\text{response\_policy} = \textit{answer}$ implies $\text{reveal\_decision} \in \{\textit{hint}, \textit{partial}, \textit{full}\}$.
A consistency loss (Section~\ref{sec:joint}) penalises violations during joint training.

\subsection{Architecture Stack}
\label{sec:architectures}

We organise our benchmark into four tracks (Figure~\ref{fig:architecture_stack}):

\begin{figure}[t]
\centering
\includegraphics[width=\linewidth]{figures/architecture_stack.pdf}
\caption{Four-track benchmark stack. The same scenario-generated data feeds architecture benchmarking, social-state encoders, response models, and the structured Qwen3 pipeline.}
\label{fig:architecture_stack}
\end{figure}

\paragraph{Track A: From-scratch SLMs (15--22M).}
Four architectures share the same vocabulary (GPT-2 tokeniser, 50,257 tokens), embedding dimension ($d=512$), and training data (16,905 dialogue lines):

\begin{enumerate}
    \item \textbf{GPT}: 6-layer transformer decoder, 8 attention heads.
    \item \textbf{PrefixGPT}: GPT + learned soft-prefix from an 8D conditioning vector (5 OCEAN + 3 VAD), projected to 4 prefix tokens via a 2-layer MLP.
    \item \textbf{MoE}: GPT with Mixture-of-Experts FFN layers (4 experts, top-2 routing with load-balancing loss).
    \item \textbf{Mamba-like}: Simplified state-space model (SSM) with selective scan, implemented in pure PyTorch.
\end{enumerate}

\paragraph{Track B: Conditioning encoders.}
DistilBERT-based encoders predict personality (OCEAN, 5D regression) and affect (VAD, 3D regression) from dialogue context.
These produce the conditioning vector $c_t \in \mathbb{R}^8$ used by PrefixGPT and ConditionalDialogue.

\paragraph{Track C: Response generation.}
Three models generate NPC responses:
\begin{enumerate}
    \item \textbf{ConditionalDialogue}: From-scratch GPT + personality/affect soft-prefix embeddings ($c_t \to$ prefix via learned projection).
    \item \textbf{TinyLlama + LoRA}: 1.1B pretrained chat model, fine-tuned with LoRA ($r=16$, $\alpha=32$) on SFT data without social-state conditioning.
    \item \textbf{Gemma-2-2B-it + QLoRA}: 2B pretrained instruction model with QLoRA fine-tuning.
    \item \textbf{Gemma-4-E2B + QLoRA}: exploratory larger active-parameter baseline.
\end{enumerate}

\paragraph{Track D: Structured LLM.}
A three-stage pipeline on Qwen3:
\begin{enumerate}
    \item \textbf{Latent predictor}: Qwen3-1.7B backbone + LoRA + 29 classification heads. Each head is a 2-layer MLP ($\text{hidden} \to 256 \to n_\text{classes}$) applied to the last-token hidden state.
    \item \textbf{Response generator}: Qwen3-1.7B + LoRA ($r=32$), supervised on input $\to$ response with gold $Z_t$ in the prompt.
    \item \textbf{Joint model}: Shared backbone with both classification heads and LM loss, plus a consistency loss $\mathcal{L}_\text{consist}$ penalising $P(\text{full\_reveal}) \cdot P(\text{high\_secrecy})$.
\end{enumerate}

\subsection{Training}
\label{sec:training}

\paragraph{Loss functions.}
Track A uses standard cross-entropy LM loss.
Track B uses MSE for regression targets.
Track D Stage~1 uses a weighted multi-head loss:
\begin{equation}
    \mathcal{L}_\text{heads} = \sum_{g \in \{C,A,M,R,N,D\}} \lambda_g \cdot \frac{1}{|g|} \sum_{f \in g} \mathcal{L}_\text{CE}(f)
\end{equation}
with $\lambda_R = \lambda_D = 2.0$, $\lambda_M = 1.5$, others $= 1.0$, reflecting the importance of relational and decision heads.

\paragraph{Joint loss (Stage 3).}
\label{sec:joint}
\begin{equation}
    \mathcal{L}_\text{joint} = \mathcal{L}_\text{heads} + \lambda_Y \mathcal{L}_\text{LM} + \lambda_\text{consist} \cdot \mathbb{E}\left[P(\text{full\_reveal}) \cdot P(\text{high\_secrecy})\right]
\end{equation}

\paragraph{Optimisation.}
All models use AdamW.
Track A: $\text{lr}=3 \times 10^{-4}$, batch 32, 20 epochs.
Track D: $\text{lr}=2 \times 10^{-4}$ (backbone), $4 \times 10^{-4}$ (heads), cosine schedule, gradient accumulation $=32$, BFloat16.

\begin{figure}[t]
\centering
\includegraphics[width=\linewidth]{figures/data_flow.pdf}
\caption{Data flow from scenario specifications to validated splits and paper artifacts. Each generated turn is validated against the $Z_t$ schema before use in training or evaluation.}
\label{fig:data_flow}
\end{figure}

\subsection{Fast/Slow Routing}
\label{sec:routing}

A rule-based router classifies turns into \textit{fast path} (direct generation) or \textit{slow path} (full $Z_t$ inference) based on predicted $D_t$ and $N_t$:
\begin{equation}
    \text{slow}(D_t, N_t) = \begin{cases}
        1 & \text{if } \text{value\_conflict} = \textit{strong} \\
        1 & \text{if } \text{secrecy} = \textit{high} \wedge \text{reveal} \in \{\textit{hint, partial, full}\} \\
        1 & \text{if } \text{policy} \in \{\textit{threaten, negotiate}\} \\
        0 & \text{otherwise}
    \end{cases}
\end{equation}

% ═══════════════════════════════════════════════════════════════════════════════
% 4. EXPERIMENTAL SETUP
% ═══════════════════════════════════════════════════════════════════════════════
\section{Experimental Setup}
\label{sec:experiments}

\paragraph{Data.}
We use 7,742 dialogue turns across 736 episodes spanning 7 scenario types (secret extraction, trust building, alliance negotiation, threat escalation, deception detection, rumor confrontation, apology repair).
Each turn is annotated with the full 29-field $Z_t$ schema.
Data is split 80/10/10 into train (6,175 turns, 587 episodes), validation (683 turns, 69 episodes), and test (884 turns, 80 episodes).
SLM training uses a separate corpus of 16,905 plain-text dialogue lines.

\paragraph{Hardware.}
All experiments run on NVIDIA A30 GPUs (24\,GB VRAM), CUDA 12.4, PyTorch 2.6.

\paragraph{Metrics.}
\textit{Track A}: validation perplexity (PPL).
\textit{Track B}: F1 (personality), CCC (affect), MSE.
\textit{Track C}: PPL, ROUGE-L, BLEU-1/2/4, Distinct-1/2/3, secret leakage rate, contradiction rate.
\textit{Track D}: per-head accuracy, macro-F1, estimated Cohen's $\kappa$, per-group means, response policy F1.
\textit{Routing}: precision, recall, F1, false positive rate.

% ═══════════════════════════════════════════════════════════════════════════════
% 5. RESULTS
% ═══════════════════════════════════════════════════════════════════════════════
\section{Results}
\label{sec:results}

\subsection{Track A: From-Scratch SLM Comparison (RQ1)}

\begin{table}[t]
\centering
\caption{From-scratch SLM validation perplexity (20 epochs, seed=42). \textbf{Bold}: best.}
\label{tab:slm}
\small
\begin{tabular}{@{}lrrrrr@{}}
\toprule
\textbf{Architecture} & \textbf{Params (M)} & \textbf{val\_loss $\downarrow$} & \textbf{val\_ppl $\downarrow$} & \textbf{$\Delta$ vs GPT} & \textbf{Conditioning} \\
\midrule
\textbf{MoE}       & 22.4 & 3.739 & \textbf{42.07} & $-7.2\%$ & none \\
PrefixGPT          & 16.6 & 3.796 & 44.54          & $-1.7\%$ & OCEAN+VAD \\
GPT                & 16.1 & 3.814 & 45.32          & ---      & none \\
Mamba-like         & 15.4 & 3.975 & 53.25          & $+17.5\%$& none \\
\bottomrule
\end{tabular}
\end{table}

\textbf{Finding (RQ1):} MoE achieves the best perplexity (42.07), a 7.2\% improvement over the standard GPT baseline (Table~\ref{tab:slm}).
Sparse expert routing benefits from the multi-character, multi-style nature of NPC dialogue---different experts can specialise in different conversational registers.
PrefixGPT's modest gain ($-1.7\%$) over GPT suggests that the 8D conditioning signal provides marginal benefit at this parameter scale.
The Mamba-like SSM converges fastest (best loss at epoch 10) but plateaus at a higher perplexity, consistent with prior observations that SSMs can underperform transformers on tasks requiring long-range attention over structured input~\cite{gu2023mamba}.

\begin{figure}[t]
\centering
\includegraphics[width=0.78\linewidth]{figures/slm_ppl_comparison.pdf}
\caption{Track A validation perplexity. MoE has the lowest PPL but also the largest parameter count among the from-scratch SLMs due to expert feed-forward blocks.}
\label{fig:slm_ppl}
\end{figure}

\subsection{Track B: Conditioning Encoders}

\begin{table}[t]
\centering
\caption{Conditioning encoder evaluation.}
\label{tab:encoders}
\small
\begin{tabular}{@{}lrrrrr@{}}
\toprule
\textbf{Encoder} & \textbf{Base} & \textbf{val\_F1 $\uparrow$} & \textbf{val\_CCC $\uparrow$} & \textbf{val\_MSE $\downarrow$} & \textbf{Best Epoch} \\
\midrule
Personality (OCEAN) & DistilBERT & 0.678 & --- & 0.248 & 4/15 \\
Affect (VAD)        & DistilBERT & ---   & 0.559 & 0.005 & 13/15 \\
\bottomrule
\end{tabular}
\end{table}

Personality prediction achieves F1~$=0.678$, consistent with known difficulty of text$\to$OCEAN regression~\cite{mairesse2007using}.
The affect encoder achieves CCC~$=0.559$, enabling reliable valence-arousal-dominance tracking.
A cache of 414 NPC personality embeddings was built from training data for efficient inference.

\subsection{Track C: Response Generation (RQ3)}

\begin{table}[t]
\centering
\caption{Response generation comparison. \textbf{Bold}: best PPL. The conditioning gain is $12.3\%$.}
\label{tab:response_ppl}
\small
\begin{tabular}{@{}llrrr@{}}
\toprule
\textbf{Model} & \textbf{Conditioning} & \textbf{val\_loss $\downarrow$} & \textbf{val\_ppl $\downarrow$} & \textbf{Epochs} \\
\midrule
\textbf{ConditionalDialogue} & OCEAN+VAD soft-prefix & 1.064 & \textbf{2.90} & 5 \\
TinyLlama 1.1B + LoRA       & none (SFT)            & 1.195 & 3.30          & 3 \\
Gemma-2-2B-it + QLoRA        & NPC profile (SFT)     & 1.854 & 6.38          & 2 \\
Gemma-4-E2B + QLoRA          & NPC profile (SFT)     & 2.787 & 16.24         & 1 \\
\bottomrule
\end{tabular}
\end{table}

\textbf{Finding (RQ3):} Explicit social-state conditioning reduces perplexity by \textbf{12.3\%} (3.30$\to$2.90; Table~\ref{tab:response_ppl}).
The ConditionalDialogue model, which projects the 8D personality+affect vector into 4 soft-prefix tokens prepended to the transformer input, consistently outperforms the unconditioned TinyLlama SFT baseline despite having far fewer parameters.
The corrected Gemma-2-2B-it run is substantially stronger than the exploratory Gemma-4-E2B run (6.38 vs.\ 16.24 PPL), but both remain behind the task-specific conditional model on this small validation split.

\begin{figure}[t]
\centering
\includegraphics[width=0.86\linewidth]{figures/response_ppl_comparison.pdf}
\caption{Track C response-model validation perplexity. The fixed Gemma-2-2B-it run should be treated as the primary Gemma baseline; the Gemma-4-E2B run is retained as an exploratory engineering result.}
\label{fig:response_ppl}
\end{figure}

\begin{figure}[t]
\centering
\includegraphics[width=\linewidth]{figures/best_model_diagram.pdf}
\caption{Best response model. OCEAN and VAD signals are compressed into an 8D social vector and projected into soft-prefix tokens before LoRA response generation.}
\label{fig:best_model}
\end{figure}

\begin{table}[t]
\centering
\caption{Response generation quality metrics (Qwen3-1.7B response model, $n=100$ sampled generations).}
\label{tab:response_quality}
\small
\begin{tabular}{@{}lr@{}}
\toprule
\textbf{Metric} & \textbf{Value} \\
\midrule
ROUGE-L           & 0.120 $\pm$ 0.058 \\
ROUGE-L 95\% CI   & [0.106, 0.133] \\
BLEU-1 / BLEU-2 / BLEU-4 & 0.177 / 0.049 / 0.009 \\
Distinct-1 / Distinct-2 / Distinct-3 & 0.128 / 0.469 / 0.706 \\
Repetition rate (3-gram / 5-gram)    & 0.667 / 0.302 \\
Avg gen length / ref length (tokens) & 118.4 / 50.7 \\
Length ratio (gen/ref)                & 2.334 \\
Secret leakage rate  & 0.000 \\
Contradiction rate   & 0.000 \\
Mean policy consistency & 0.461 \\
\bottomrule
\end{tabular}
\end{table}

Table~\ref{tab:response_quality} shows generation quality for the Qwen3 response model.
ROUGE-L is low (0.120), expected for open-ended dialogue where multiple valid responses exist.
Distinct-2/3 scores (0.469/0.706) indicate reasonable lexical diversity, though the high 3-gram repetition rate (0.667) and length ratio (2.33$\times$) suggest the model tends to generate verbose, repetitive outputs---a known issue with small LMs at early training stages.
Critically, \textbf{zero secret leakage} and \textbf{zero contradictions} are maintained.

\subsection{Track D: Latent State Prediction (RQ2)}

\begin{table}[t]
\centering
\caption{Per-group latent state prediction metrics (Qwen3-1.7B, evaluated over 28 single-label heads; dialogue act is multi-label and excluded from this table).}
\label{tab:latent_groups}
\small
\begin{tabular}{@{}llrrrr@{}}
\toprule
\textbf{Group} & \textbf{Description} & \textbf{\#Heads} & \textbf{Mean Acc $\uparrow$} & \textbf{Std} & \textbf{Mean $\kappa$ $\uparrow$} \\
\midrule
$C_t$ & Conversational   & 2  & 0.767 & 0.124 & 0.718 \\
$A_t$ & Affect           & 4  & 0.778 & 0.060 & 0.667 \\
$M_t$ & Mental model     & 3  & 0.725 & 0.059 & 0.641 \\
$R_t$ & Relational stance& 12 & 0.640 & 0.068 & 0.550 \\
$N_t$ & Normative        & 4  & 0.782 & 0.087 & 0.674 \\
$D_t$ & Decision policy  & 3  & 0.673 & 0.057 & 0.598 \\
\midrule
\multicolumn{2}{l}{\textbf{Overall (28 heads)}} & 28 & \textbf{0.702} & 0.096 & \textbf{0.611} \\
\bottomrule
\end{tabular}
\end{table}

\textbf{Finding (RQ2):} The 29-head predictor achieves mean accuracy $0.702$ and mean estimated $\kappa = 0.611$ (Table~\ref{tab:latent_groups}).
The strongest groups are normative pressures ($N_t$: $\kappa=0.674$) and affect ($A_t$: $\kappa=0.667$), while relational stance ($R_t$: $\kappa=0.550$) is hardest---consistent with the difficulty of inferring interpersonal dynamics from text alone.

The top-performing individual heads are \texttt{risk\_type} (accuracy 0.892, $\kappa=0.865$) and \texttt{face\_pressure} (accuracy 0.868, $\kappa=0.802$), while the most challenging heads are \texttt{familiarity\_delta} (accuracy 0.499, $\kappa=0.374$) and \texttt{dominance\_delta} (accuracy 0.572, $\kappa=0.466$).
Response policy---the most actionable head (10 classes)---achieves accuracy 0.707, well above the chance baseline of $0.10$.
Response-policy macro-F1 differs across artifacts: the checkpoint epoch logs peak at 0.448, while the independent flat evaluation file reports 0.512.
We therefore treat response-policy F1 as approximately 0.45--0.51 until a clean checkpoint rerun reconciles the logs.

The overfitting gap (train loss 1.68 vs.\ val loss 7.39 at epoch 3) indicates the model has capacity to memorise the training distribution; scaling to a larger dataset or using the 4B backbone would likely improve generalisation.

\begin{figure}[t]
\centering
\includegraphics[width=0.86\linewidth]{figures/latent_group_scores.pdf}
\caption{Latent prediction by social-state group. Relational stance is the hardest group, while affect and normative-pressure heads are more predictable from local context.}
\label{fig:latent_groups}
\end{figure}

\begin{figure}[t]
\centering
\includegraphics[width=0.90\linewidth]{figures/qwen_latent_training.pdf}
\caption{Qwen3 latent-predictor training dynamics. Mean accuracy keeps improving through epoch 4, while response-policy macro-F1 peaks earlier, suggesting a trade-off between broad head accuracy and decision-policy calibration.}
\label{fig:qwen_training}
\end{figure}

\subsection{Routing}

The rule-based fast/slow router achieves perfect precision and recall ($\text{F1}=1.0$) on the validation set, with a slow-path rate of 53.7\%.
This is expected since the router applies deterministic rules to the same latent state used for ground truth construction.
In deployment, the router would operate on \textit{predicted} $Z_t$, where routing errors would be bounded by the latent predictor's accuracy.

\begin{figure}[t]
\centering
\includegraphics[width=\linewidth]{figures/structured_llm_pipeline.pdf}
\caption{Structured Qwen3 pipeline. The joint checkpoint is trained, but this paper reports independent latent and response evaluations rather than claiming joint-model quality improvements.}
\label{fig:structured_llm}
\end{figure}

% ═══════════════════════════════════════════════════════════════════════════════
% 6. DISCUSSION
% ═══════════════════════════════════════════════════════════════════════════════
\section{Discussion}
\label{sec:discussion}

\paragraph{Architecture choice at small scale.}
The MoE advantage (7.2\% over GPT) at 15--22M parameters suggests that expert specialisation captures the multi-register nature of NPC dialogue (e.g., different NPCs have distinct speaking styles).
However, the gap is modest compared to scaling the backbone (from-scratch 42.07 PPL vs.\ pretrained ConditionalDialogue 2.90 PPL), confirming that architecture choice matters less than pretraining data at small scales.

\paragraph{Conditioning effectiveness.}
The 12.3\% perplexity reduction from explicit OCEAN+VAD conditioning is notable because both models (ConditionalDialogue and TinyLlama SFT) have access to the same training data.
The conditioning signal provides a compressed social-state summary that the model can use as a prior, reducing uncertainty about the NPC's likely response style.

\paragraph{Latent state as an interpretable bottleneck.}
The 29-head predictor demonstrates that structured social state is learnable from dialogue context.
Mean estimated $\kappa=0.611$ indicates substantial lift over a uniform chance baseline, suggesting the latent state captures genuine social dynamics rather than noise.
The interpretability of individual heads (e.g., ``\texttt{secrecy\_pressure=high}'' or ``\texttt{response\_policy=withhold}'') enables human auditing of NPC behaviour.

\paragraph{Secret-keeping.}
Zero keyword-detected secret leakage across sampled generations supports the design of our reveal\_decision head and consistency constraints, but does not constitute a formal leakage guarantee.
In a deployed system, the latent predictor's $D_t$ output would gate information disclosure, providing a verifiable safety mechanism.

\paragraph{Limitations.}
\begin{enumerate}
    \item \textbf{Data scale}: 7,742 turns is small by NLP standards; results may shift with $10\times$ data.
    \item \textbf{Single seed}: Track A results are for seed=42 only; multi-seed experiments would strengthen claims.
    \item \textbf{Joint-model evaluation}: The three-stage Qwen3 joint model has been trained, but it has not yet been independently evaluated for generation quality or per-head latent accuracy.
    \item \textbf{Synthetic labels}: Social-state annotations are LLM-generated (not human-annotated), introducing label noise.
    \item \textbf{Evaluation coverage}: ROUGE-L and BLEU are weak proxies for dialogue quality; human preference and task-success evaluation are needed for deployment-ready claims.
    \item \textbf{Repetition}: The response model shows high 3-gram repetition (66.7\%) and generates responses 2.3$\times$ longer than references, indicating a need for length penalties or nucleus sampling tuning.
\end{enumerate}

% ═══════════════════════════════════════════════════════════════════════════════
% 7. CONCLUSION
% ═══════════════════════════════════════════════════════════════════════════════
\section{Conclusion}
\label{sec:conclusion}

We benchmarked 12 models across four tracks for social-state-conditioned NPC dialogue generation:
\begin{enumerate}
    \item \textbf{MoE achieves best from-scratch perplexity} (42.07, $-7.2\%$ vs.\ GPT) at 15--22M parameters.
    \item \textbf{Explicit social-state conditioning reduces perplexity by 12.3\%} (3.30$\to$2.90) via OCEAN+VAD soft-prefix.
    \item \textbf{Structured social state is learnable}: mean accuracy $0.702$, mean estimated $\kappa = 0.611$ across evaluated heads at the 1.7B scale.
    \item \textbf{Zero keyword-detected secret leakage} is observed in sampled generations, supporting the reveal\_decision mechanism.
\end{enumerate}

The structured latent state $Z_t$ provides an interpretable, auditable intermediate representation that bridges perception and generation.
Future work will: (a)~evaluate the joint model end-to-end, (b)~scale to the Qwen3-4B backbone, (c)~incorporate human evaluation, and (d)~deploy in an interactive game environment.

% ═══════════════════════════════════════════════════════════════════════════════
% REFERENCES
% ═══════════════════════════════════════════════════════════════════════════════
\bibliography{references}
\bibliographystyle{plainnat}

% ═══════════════════════════════════════════════════════════════════════════════
% APPENDIX
% ═══════════════════════════════════════════════════════════════════════════════
\appendix

\section{Per-Head Evaluation Details}
\label{app:perhead}

\begin{figure}[h]
\centering
\includegraphics[width=0.88\linewidth]{figures/latent_head_heatmap.pdf}
\caption{Per-head latent accuracy and estimated $\kappa$. Relational-delta heads are consistently harder than affect and normative heads.}
\label{fig:latent_heatmap}
\end{figure}

\begin{table}[h]
\centering
\caption{Per-head latent state prediction metrics, sorted by accuracy (descending). $\kappa$ is estimated from accuracy and uniform chance baseline.}
\label{tab:perhead}
\scriptsize
\begin{tabular}{@{}lrrrrr@{}}
\toprule
\textbf{Head} & \textbf{Group} & \textbf{\#Cls} & \textbf{Acc $\uparrow$} & \textbf{Chance} & \textbf{$\kappa$ (est.)} \\
\midrule
risk\_type        & $C_t$ & 5  & 0.892 & 0.200 & 0.865 \\
face\_pressure    & $N_t$ & 3  & 0.868 & 0.333 & 0.802 \\
arousal           & $A_t$ & 3  & 0.854 & 0.333 & 0.780 \\
duty\_pressure    & $N_t$ & 3  & 0.823 & 0.333 & 0.734 \\
valence           & $A_t$ & 3  & 0.810 & 0.333 & 0.715 \\
value\_conflict   & $N_t$ & 3  & 0.801 & 0.333 & 0.701 \\
player\_credibility & $M_t$ & 3 & 0.799 & 0.333 & 0.699 \\
threat            & $A_t$ & 3  & 0.754 & 0.333 & 0.631 \\
respect\_delta    & $R_t$ & 5  & 0.727 & 0.200 & 0.659 \\
reveal\_decision  & $D_t$ & 4  & 0.720 & 0.250 & 0.627 \\
player\_intent    & $M_t$ & 9  & 0.719 & 0.111 & 0.684 \\
affection\_level  & $R_t$ & 5  & 0.716 & 0.200 & 0.645 \\
response\_policy  & $D_t$ & 10 & 0.707 & 0.100 & 0.675 \\
obligation\_delta & $R_t$ & 5  & 0.707 & 0.200 & 0.634 \\
affection\_delta  & $R_t$ & 5  & 0.696 & 0.200 & 0.620 \\
control           & $A_t$ & 3  & 0.694 & 0.333 & 0.541 \\
trust\_level      & $R_t$ & 5  & 0.683 & 0.200 & 0.603 \\
dominance\_level  & $R_t$ & 5  & 0.671 & 0.200 & 0.589 \\
player\_knowledge & $M_t$ & 4  & 0.656 & 0.250 & 0.541 \\
tone              & $C_t$ & 6  & 0.643 & 0.167 & 0.571 \\
secrecy\_pressure & $N_t$ & 3  & 0.638 & 0.333 & 0.456 \\
respect\_level    & $R_t$ & 5  & 0.630 & 0.200 & 0.538 \\
familiarity\_level& $R_t$ & 5  & 0.598 & 0.200 & 0.497 \\
repair\_strategy  & $D_t$ & 5  & 0.593 & 0.200 & 0.491 \\
trust\_delta      & $R_t$ & 5  & 0.592 & 0.200 & 0.490 \\
obligation\_level & $R_t$ & 5  & 0.586 & 0.200 & 0.482 \\
dominance\_delta  & $R_t$ & 5  & 0.572 & 0.200 & 0.466 \\
familiarity\_delta& $R_t$ & 5  & 0.499 & 0.200 & 0.374 \\
\bottomrule
\end{tabular}
\end{table}

\section{Sample Generations}
\label{app:samples}

Table~\ref{tab:samples} shows representative generated responses from the Qwen3 response model.

\begin{table}[h]
\centering
\caption{Sample NPC responses. Latent state context is abbreviated.}
\label{tab:samples}
\scriptsize
\begin{tabular}{@{}p{3.5cm}p{4cm}p{4cm}r@{}}
\toprule
\textbf{Latent Context} & \textbf{Gold Response} & \textbf{Generated Response} & \textbf{ROUGE-L} \\
\midrule
policy=negotiate, reveal=none, secrecy=medium &
``I can mend your blade, but I won't do it for nothing. You'll have to trade something of value...'' &
``Okay, I can fix your sword, but I need a fair price for my time and materials. How much do you think you can offer?...'' &
0.103 \\
\midrule
policy=clarify, reveal=hint, secrecy=high &
``I've heard whispers, not from the council, but from those who know the iron in their bones...'' &
``The Iron Legion is a band of warriors from the east... I've heard whispers that they're looking for a blade of ancient power...'' &
0.104 \\
\bottomrule
\end{tabular}
\end{table}

\section{Data Generation Pipeline}
\label{app:datagen}

Dialogue data is generated using a scenario-based pipeline with 7 scenario types:
secret extraction, trust building, alliance negotiation, threat escalation, deception detection, rumor confrontation, and apology repair.
Each scenario defines the NPC's profile (goals, values, secrets, persona), the initial relationship stance, and the player's conversational strategy.
An LLM generates multi-turn dialogues with per-turn $Z_t$ annotations, which are validated against the schema's consistency rules.
Counterfactual augmentation flips single dimensions (e.g., reversing secrecy pressure) to diversify the training distribution.

\end{document}
